\section{Method}
\label{sec:method}

All measurements come from one machine, an Apple M3 Pro whose performance
cores run at 4.0\,GHz. A small pointer-chasing program measured what a load
costs on it: about 4 cycles from the L1 cache (128\,KB), 15--36 cycles from L2
(16\,MB) and 360--450 cycles from main memory, but only 4--8 cycles at any size
when memory is read in order, because the processor then fetches the next data
in advance. These numbers are the scale for every cost in this report.

Timings use the optimized build of the solver (\texttt{-O2}, assertions off),
and profiles use the same build with debug symbols added. Instead of seconds
I report processor cycles and instructions, read from \texttt{/usr/bin/time},
since cycles do not depend on the clock speed, and I keep the fastest of three
runs. The solver also counts its own operations. These counts are identical on
every run of the same input, so whenever a change should not affect the search
(proof output, memory padding, the counters themselves) I checked that they
stayed identical. The counters cost 3--4\% of the run time; my first version,
which updated memory on every clause visit, cost 10\%.

The solver is deterministic, but the same formula written differently
(variables renamed, signs flipped, clauses reordered) can need twice as many
conflicts, because the heuristics take a different path. I therefore compare
configurations on shuffled copies of an instance, generated from fixed seeds,
and in pairs: both configurations run on the same copy, so each copy's own
difficulty cancels out in their ratio. A sign test and a Wilcoxon signed-rank
test on the log ratios tell whether an effect is real, and a bootstrap 95\%
confidence interval for the median ratio gives its size.

Every experiment is a script that writes a CSV file recording the git commit,
build flags, machine and command, and every figure, table and number in this
report is generated from those files (Appendix~\ref{app:repro}). Three limits
apply throughout. The processor's hardware event counters were not available
(the installed Xcode was too old for the operating system), so statements
about caches and branches are inferences from several independent
measurements, not direct counts. There is a single machine. And all instances
have at most 250 variables, so most of them fit in the L2 cache.
